\begin{pastexam}[label={px:ntu-111-dsaa-10}]{臺大資工所 111 學年度 第 10 題}
（單選）

(5 points) Consider an integer array $A$ of size $n>0$ with elements $A[1],\ldots,A[n]$, where
the first $p$ elements are positive and the other $(n-p)$ elements are negative. The following
algorithm calculates $p$ correctly by calling \textsc{Compute-P}$(A,n)$. What loop invariant is
maintained for the \textbf{while} loop of the algorithm?

\begin{quote}
\textsc{Compute-P}$(A, len)$\\
1\quad \textbf{if} $A[1]<0$ \textbf{return} 0\\
2\quad $left=1, right=len$\\
3\quad \textbf{while} $left<right$\\
4\quad\quad $m=\text{ceiling}(\frac{left+right}{2})$\\
5\quad\quad \textbf{if} $A[m]>0$\\
6\quad\quad\quad $left=m$\\
7\quad\quad \textbf{else}\\
8\quad\quad\quad $right=m-1$\\
9\quad \textbf{return} $left$
\end{quote}
\begin{enumerate}[label=(\Alph*)]
\item $left+1\le p\le right$
\item $left+1\le p\le right-1$
\item $left\le p\le right$
\item $left\le p\le right-1$
\item none of the other choices
\end{enumerate}
\end{pastexam}

\begin{pastsoln}
答案是 (C)。

第 1 行保證進入迴圈時 $A[1] > 0$，所以 $p \ge 1$。 正數恰好是前 $p$ 個，因此 $A[m] > 0$ 等價於
$m \le p$，$A[m] < 0$ 等價於 $m \ge p + 1$（題目說每個元素不是正就是負，沒有 $0$）。 照資結書
第~1.4~節的三步驗證 (C)：
\begin{itemize}
\item \textbf{初始}：$left = 1 \le p$，$p \le n = right$。
\item \textbf{維持}：迴圈裡 $left < right$，$m = \lceil (left + right)/2 \rceil$ 滿足
$left < m \le right$。 若 $A[m] > 0$，則 $m \le p$，新的 $left = m \le p$，$right$ 不變；若
$A[m] < 0$，則 $m \ge p + 1$，新的 $right = m - 1 \ge p$，$left$ 不變。 兩種情形
$left \le p \le right$ 都仍成立。
\item \textbf{結束}：每一輪 $left$ 嚴格變大（$m > left$）或 $right$ 嚴格變小，迴圈一定會停，停在
$left = right$，代入不變量得 $p = left$，第 9 行回傳正確。
\end{itemize}

其他選項：
\begin{itemize}
\item (A) $left + 1 \le p$：$p = 1$ 時（例如 $A = \langle 5, -2 \rangle$），一開始 $left = 1$，
$left + 1 = 2 > p$，第一次檢查就不成立。 而且迴圈結束時 $left = p$，這個不等式一定不成立。
\item (B) 比 (A) 多要求 $p \le right - 1$，同樣在 $p = 1$ 時一開始就不成立；結束時
$left = right$，$left + 1 \le right - 1$ 不可能成立。
\item (D) $p \le right - 1$：全部都是正數時（例如 $A = \langle 3, 1 \rangle$，$p = n = 2$），一開始
$right = 2$，$p \le 1$ 不成立。 結束時 $left = right$，$left \le p \le left - 1$ 也不可能成立。
\item (E)：(C) 成立，所以不選。
\end{itemize}

這個程序找的是「最後一個滿足條件（正數）的位置」：往右縮寫成 $left = m$，所以中點一定要取天花板
（資結書第~17.1.3~節）。 如果第 4 行改成地板，$left = 1$、$right = 2$、$A[1] > 0$ 時 $m = 1$，
$left$ 不動，迴圈不會停。
\end{pastsoln}
